\documentclass[11pt]{article}
\usepackage[T1]{fontenc}
\usepackage[utf8]{inputenc}
\usepackage{lmodern,microtype}
\usepackage[a4paper,margin=28mm]{geometry}
\usepackage{amsmath,amssymb,amsthm,mathtools,booktabs,longtable,array}
\usepackage[dvipsnames]{xcolor}
\usepackage{xr-hyper}
\usepackage[colorlinks=true,linkcolor=MidnightBlue,citecolor=MidnightBlue,urlcolor=MidnightBlue]{hyperref}
\newtheorem{theorem}{Theorem}[section]
\newtheorem{proposition}[theorem]{Proposition}
\newtheorem{lemma}[theorem]{Lemma}
\newtheorem{corollary}[theorem]{Corollary}
\theoremstyle{definition}\newtheorem{definition}[theorem]{Definition}
\theoremstyle{remark}\newtheorem{remark}[theorem]{Remark}
\newtheorem{question}[theorem]{Question}
\newtheorem{conjecture}[theorem]{Conjecture}
\DeclareMathOperator{\FPdim}{FPdim}
\DeclareMathOperator{\Irr}{Irr}
\DeclareMathOperator{\Tr}{Tr}
\DeclareMathOperator{\Rep}{Rep}
\let\Vec\relax
\DeclareMathOperator{\Vec}{Vec}
\newcommand{\C}{\mathcal C}
\newcommand{\Z}{\mathcal Z}
\newcommand{\Q}{\mathbb Q}
\newcommand{\N}{\mathbb Z_{\geq0}}
\newcommand{\R}{\mathbb R}
\newcommand{\CC}{\mathbb C}
\newcommand{\ZZ}{\mathbb Z}
\newcommand{\id}{\operatorname{id}}
\allowdisplaybreaks[2]
\setlength{\emergencystretch}{2em}

\hypersetup{pdftitle={Rank-five Grothendieck rings through multiplicity twenty: two general residuals at eighteen}}
\title{Rank-five Grothendieck rings through multiplicity twenty:\\two general residuals at eighteen}
\date{23 September 2026}
\newcommand{\F}[2]{F^{(5)}_{#1,#2}}
\newcommand{\FF}{\mathbb F}
\newcommand{\OO}{\mathcal O}
\newcommand{\supp}{\operatorname{supp}}
\begin{document}
\maketitle
\begin{abstract}
We classify the rank-five fusion rings of multiplicity at most seventeen
that admit categorifications over $\CC$, without assuming pivotality.
Exactly seventeen of the $23413$ rings admit categorifications, and each
has a unitary realization. The exclusions are automatic applications of
arithmetic and central-induction criteria. Two new reductions make the
method effective: a mod-two spectral bound and parity rigidity force
sphericality for a class of quadratic fusion rings, while a central-support
analysis for cyclic order-four near-group rings supplies finite indicator
systems. A single rational/quadratic relative-trace construction contains
the higher-indicator obstructions. No exclusion is assigned by a census
label or imported from a classification list. The pinned
\texttt{fusion-ring-census} release records $40999$ rings through
multiplicity twenty. At exact multiplicity nineteen and twenty, every
ring is excluded from general categorification by an exact certificate.
Only two general categorification questions, both at multiplicity
eighteen, remain undecided by the implemented criteria; both rings are
excluded pivotally. Exact data, certificates, verifiers and complete
reproduction instructions accompany the paper.
\end{abstract}
\setcounter{tocdepth}{1}
\tableofcontents

\section{Introduction and results}
A fusion ring is a based ring with nonnegative integral structure
constants $N_{ij}^k$, distinguished unit and involutive duality, satisfying
Frobenius reciprocity. Its rank is the number of basis elements, and its
multiplicity is $\mu(R)=\max_{i,j,k}N_{ij}^k$. Unit coefficients are included,
so $\mu(R)\ge1$. By a \emph{Grothendieck ring} we mean one admitting a
fusion category over $\CC$. The question here is existence for a based
ring, not classification of all associators realizing it.

Small-rank enumeration provides candidates rather than categorifications.
The pinned census \cite{Census} and the multiplicity-one analysis of \cite{LPR}
make this distinction concrete. The rank-four methods of \cite{DP}
combine formal codegrees, sphericalization and central induction to
obtain general, not merely pivotal, exclusions. In rank five the main
additional issue is to avoid introducing a new argument for each surviving
ring. We address it with parameterized necessary conditions and exact
certificates whose hypotheses are read from the multiplication tensor.

\begin{theorem}\label{thm:main}
Up to based-ring isomorphism, there are $23413$ rank-five fusion rings
with $\mu\le17$. Exactly the seventeen rings in
Table~\ref{tab:models} admit fusion-category categorifications over $\CC$.
All seventeen have unitary realizations. Therefore general, pivotal,
spherical, pseudo-unitary and unitary existence give the same
classification of rings in this range.

There are $28493$ rank-five fusion rings with $\mu\le18$.
All except these seventeen and the two rings
\[
\F{18}{5073},\qquad \F{18}{5076}
\]
are excluded from general categorification by the supplied criteria.
Both displayed rings are excluded spherically and pivotally. Hence the
same seventeen rings are the complete positive list through multiplicity
twenty for pivotal, spherical, pseudo-unitary and unitary existence.

There are $5640$ and $6866$ rank-five rings of exact multiplicities
nineteen and twenty, respectively. Every one has an exact general
noncategorifiability certificate in the accompanying supplement. Thus
through multiplicity twenty, of the $40999$ census rings, exactly
seventeen are categorifiable, exactly two remain undecided for general
categorification by these criteria, and the other $40980$ are excluded.
\end{theorem}

The numbering convention is fixed in Section~\ref{sec:automatic}; at
multiplicity eighteen the two subscripts are also their independent
lexicographic indices. Their full tensors, dimensions and certificates
are supplied. ``Undecided'' means that no general exclusion certificate
was obtained by the present criteria and bounded searches. In particular,
a timed-out signed search is not treated as an exclusion, nor is it a
proof of feasibility. We make no assertion that a literature-wide
nonexistence result cannot decide either ring.

The new theory has two roles. Theorem~\ref{thm:parity} gives a uniform
sphericality criterion with only small integer and mod-two calculations;
it eliminates the need for large signed-table searches on its domain.
Theorem~\ref{thm:cyclic} turns the central support of $K(C_4,4k)$ into a
finite obstruction compatible with the same relative-trace framework.
The general finite-indicator criterion is Theorem~\ref{thm:finite}.
The rank-four obstructions used within it or within sphericalization are
stated with references, not reproved. The classification is completed by
a single exhaustive application and by tensor matches to positive
constructions, rather than by separate negative case proofs.

\section{Census input and commutativity}\label{sec:data}
The input is the pinned \texttt{fusion-ring-census} release
\cite{Census}. It counts based-ring isomorphism classes, rather than
raw multiplication tables; its exact-multiplicity counts through twenty
are displayed in Table~\ref{tab:census}. Our exclusion ledger addresses
each rank-five tensor through eighteen. To connect this archived ledger
with the repository tables, we canonically flatten every tensor under
all $24$ permutations of the four nonunit basis elements. The resulting
sets agree at each multiplicity $1\le\mu\le18$: all $28493$ repository
tensors have an isomorphic archived tensor, and conversely. The
multiplicity-nineteen and -twenty certificates address the repository
tensors directly, in their recorded repository order. Labels do
not generally agree across these releases, so the classification
transport uses based-ring isomorphism, never equality of row indices
or of labeled-tensor fingerprints. The census identifies the input;
the categorical conditions and certificates decide its members.

The following argument supplies the commutativity needed for both the
ordinary and signed sphericalization computations. The unsigned assertion
also appears in \cite[Proposition 9.2]{ABDP}.

\begin{lemma}\label{lem:commutativity}
Every rank-five fusion ring is commutative. A rank-five signed Frobenius
based ring admitting a nonzero, dual-equal dimension character is also
commutative.
\end{lemma}
\begin{proof}
The complex algebra of a (possibly signed) Frobenius based ring is
semisimple: the coefficient-of-unit functional gives the positive
inner product making the basis orthonormal and multiplication a
$*$-representation. A noncommutative algebra of dimension five with a
one-dimensional representation must be $\CC\oplus M_2(\CC)$.
Its unique one-dimensional character is Galois invariant, hence rational,
and its values are algebraic integers, hence integers.
For a fusion ring this character is FP; in the signed assertion it is
the given dimension character. In either case its formal codegree is an
integer $D\geq5$, the sum of five nonzero integer squares.
If $f$ denotes the formal codegree of the two-dimensional representation,
the standard identity $\sum_E\dim(E)/f_E=1$ gives
\[
 f=\frac{2D}{D-1}=2+\frac2{D-1}.
\]
Formal codegrees are algebraic integers. The displayed rational number
is not an integer for $D\geq5$, a contradiction.
\end{proof}

In the self-dual case, our tensor coordinate convention is the twenty
ordered nonunit triples
\[
111,112,113,114,122,123,124,133,134,144,
222,223,224,233,234,244,333,334,344,444.
\]
The coefficients $v_j$ are listed in this order; multiplication with the
unit is fixed by the based-ring axioms. The two tensors resolved by
inert-prime parity in
Section~\ref{sec:automatic} are recorded in these coordinates. For
nonself-dual inputs, the data contain the full multiplication tensors.

\begin{table}[ht]\centering\small
\begin{tabular}{rrrr}\toprule
$\mu$&Number of rings&$\mu$&Number of rings\\\midrule
1&16&11&1383\\
2&37&12&1948\\
3&82&13&2211\\
4&134&14&2554\\
5&209&15&2936\\
6&336&16&4157\\
7&477&17&4255\\
8&733&18&5080\\
9&863&19&5640\\
10&1082&20&6866\\
\bottomrule\end{tabular}
\caption{Rank-five based-isomorphism classes at exact multiplicity, from the pinned census \cite{Census}. The categorification decisions in this paper cover every displayed multiplicity.}\label{tab:census}
\end{table}


\section{Categorical inputs and a common certificate framework}
\label{sec:inputs}
Let $N_i=(N_{ij}^k)_{j,k}$ and
\[
 B=\sum_iN_iN_{i^*},\qquad p_R(t)=\det(tI-B).
\]
For a commutative fusion ring the eigenvalues of $B$, counted with
multiplicity, are its formal codegrees. The same convention will be used
for a commutative signed Frobenius quotient. We write $d_i=\FPdim(X_i)$,
$D=\sum_i d_i^2$, and use $F:\Z(\C)\to\C$ and $I$ for forgetful and
induction whenever $\C$ is spherical.

\subsection{Arithmetic and sphericalization}
We use the following necessary conditions from
\cite[Section 2]{DP}. Formal codegrees of a complex fusion category are
$d$-numbers. For a monic minimal polynomial
$t^r+a_1t^{r-1}+\cdots+a_r$, this is the integer divisibility condition
$a_r^j\mid a_j^r$ for $1\le j\le r$.
Fusion eigenvalues are cyclotomic integers. For a based fusion subring
$S\subset R$, the quotient $\FPdim(R)/\FPdim(S)$ is an algebraic integer.
The normalized-norm and finite-prime field tests of
\cite[Lemmas 2.3--2.4]{DP} give exact certificates of the required failures.
These tests depend on polynomials, divisibility and factorizations over
finite fields, not numerical approximations to roots.

Weak integrality permits the pseudo-unitary reduction. More generally,
a ring with exactly two distinct simple Frobenius--Perron dimensions
has a pseudo-unitary Galois conjugate whenever it has a categorification;
we use \cite[Proposition 2.1]{DP}. In the spherical case a dimension
character in the FP Galois orbit gives the same reduction by
\cite[Lemma 2.2]{DP}. Every dimension character is nonzero on simples
and equal on dual pairs.

For an arbitrary categorification, pass to the canonical sphericalization
and quotient its ring by the relation $\varepsilon=-1$, where
$\dim\varepsilon=-1$. The resulting signed coefficients satisfy
\begin{equation}\label{eq:signed}
 |T_{ij}^k|\le N_{ij}^k,\qquad T_{ij}^k\equiv N_{ij}^k\pmod2.
\end{equation}
The selected lifts may be chosen to have second indicator $+1$ whenever
the underlying simple is self-dual. In that normalization the adjoint
parity condition becomes
\begin{equation}\label{eq:positive-lifts}
 T_{jj^*}^i\equiv N_{jj^*}^i\pmod4
 \quad(i=i^*).
\end{equation}
This is \cite[Lemma 7.1 and Corollary 7.2]{DP}. We impose no additional
sign gauge on the odd coefficients: its compatibility with this
normalization would need separate justification. Dual pairs are selected
compatibly. The rank-five commutativity argument of Section~\ref{sec:data}
applies because a spherical dimension character is nonzero on every
simple. Thus commutative signed-table enumeration is exhaustive here.

The spherical and signed character tests are exactly those of
\cite[Sections 2, 6--7]{DP}: nonvanishing, equality on dual pairs,
formal-codegree norm and field conditions, and the positive central-square
moment inequalities, including the central sign object when necessary.
A signed table that is a sign gauge of $R$ splits the sphericalization,
so the spherical exclusion of $R$ applies. Every other signed character
component must receive its own necessary-condition certificate. This
is a quantified test over all signed tables and all components, not a
conjecture that such tables suffice for associators.

\subsection{Low induction and exact moments}
In a pseudo-unitary categorification all central dimensions are positive.
Since the ring is commutative, the induction of the unit is
multiplicity-free. Its simple dimensions are $D/f$, one for each formal
codegree $f$, and all its twists equal one
\cite[Theorem 2.13]{Ostrik}. If $L$ is the matrix of their restrictions,
including the unit column, then
\begin{equation}\label{eq:low}
 L_{0a}=1,\quad \sum_iL_{ia}d_i=D/f_a,\quad
 (LL^{\mathsf T})_{0i}=B_{0i},\quad
 Q=B-LL^{\mathsf T}\ge0\ \text{entrywise},\quad Q\succeq0.
\end{equation}
Column bounds $L_{ia}\le B_{0i}$ and $L_{ia}^2\le B_{ii}$ make this a
finite enumeration. Equal formal-codegree columns may be put in
lexicographic order to remove only redundant permutations.
The low mass in $I(X_i)$ is
$\ell_i=\sum_j(LL^{\mathsf T})_{ij}d_j$.
The moment bounds of \cite[Theorem 3.1 and Proposition 3.2]{DP} include
\begin{equation}\label{eq:low-moment}
 (3d_i+1)\ell_i\le Dd_i(d_i+1)\quad(i=i^*),
 \qquad 3\ell_i\le Dd_i\quad(i\ne i^*).
\end{equation}
Their low-column implementation is \cite[Proposition 3.3]{DP}.
We retain them as cheap consequences of the common moment framework,
not as separate ring-specific methods.

For every positive integer $j$, the Ng--Schauenburg formula is
\begin{equation}\label{eq:NS}
 \sum_{Z\in\Irr\Z(\C)}[F(Z):X_i]\dim(Z)\theta_Z^j
   =D\nu_j(X_i).
\end{equation}
Here $\nu_1(X_i)=0$ for $i\ne0$; $\nu_2(X_i)=0$ for a nonself-dual
simple and is $\pm1$ otherwise. If a self-dual simple occurs oddly in
some $Y\otimes Y^*$, its second indicator is $+1$ by the adjoint-parity
result just cited. Formula~\eqref{eq:NS} is valid for spherical fusion
categories \cite[Theorem 4.1]{NS}; the nonnegative weights used below
require the stated pseudo-unitary reduction.

The first tests in this section are inexpensive filters. Their
consequences are not claimed to be logically independent or a minimal
axiomatization. In particular, we retain an arithmetic rejection when
it avoids constructing a much larger, logically stronger moment system.

\section{Finite relative-trace certificates}\label{sec:finite}
The same construction will be used for integral and quadratic dimensions.
In the quadratic case put $K=\Q(\sqrt q)$, with $q>1$ squarefree, and let
$F=q$ for $q\equiv1\pmod4$, and $F=4q$ otherwise. In the rational case
put $K=\Q$ and $F=1$. A rational computation is not interpreted as a
quadratic field with two embeddings.

\subsection{A finite set of joint traces}
Let $P$ be a finite set of rational primes containing those dividing $F$,
and fix a height $h\ge3$. For a root of unity $\zeta$ with order supported
on $P$, let
\[
 \tau_j(\zeta)=\frac{\Tr_{L/K}(\zeta^j)}{[L:K]},\qquad 1\le j\le h,
\]
where $L$ is any common cyclotomic field containing $K$ and $\zeta$.
The value is independent of enlarging $L$.

\begin{lemma}\label{lem:trace-box}
All these joint traces are contained in an explicitly computable finite
set. It suffices to use orders dividing
\begin{equation}\label{eq:trace-box}
 B_h=\prod_{p\in P}p^{\max(v_p(F),1)+\lfloor\log_p h\rfloor}
\end{equation}
and to add the all-zero vector. For $K=\Q$, the trace of a primitive
$n$th root is $\mu(n)/\varphi(n)$. For quadratic $K$, if $\chi$ is the
primitive quadratic character of conductor $F$, then
\begin{equation}\label{eq:gauss}
 \tau_1(\zeta_n^a)=\frac{\mu(n)}{\varphi(n)}+
 \begin{cases}
 \displaystyle\frac{\chi(a)\mu(g)\chi(g)\sqrt F}{\varphi(n)},
       &n=Fg,\ (F,g)=1,\\[3pt]
 0,&\text{otherwise},
 \end{cases}
\end{equation}
for $(a,n)=1$. Higher coordinates are obtained by reducing
$\zeta_n^{aj}$ to its exact order. The same sign $\chi(a)$ is used in
every nonrational coordinate of a joint trace.
\end{lemma}
\begin{proof}
For $F\nmid n$, the intersection $K\cap\Q(\zeta_n)$ is $\Q$, since the
conductor of $K$ is $F$, so the normalized trace is the normalized
Ramanujan sum. For $F\mid n$, average over the subgroup on which $\chi$
is one. Its indicator is $(1+\chi)/2$. The untwisted sum is $\mu(n)$,
and the induced quadratic Gauss sum is
$\chi(a)\mu(g)\chi(g)\sqrt F$ when $(F,g)=1$ and zero otherwise.
This gives \eqref{eq:gauss}, including nonsquarefree $g$.
For a power $j$, set $r=(n,j)$ and apply the formula at order $n/r$,
with primitive exponent $aj/r$; this proves the simultaneous sign rule.

Suppose $v_p(n)>\max(v_p(F),1)+\lfloor\log_p h\rfloor$ for some $p$.
For every $j\le h$, its reduced order still has $p$-valuation strictly
greater than both $v_p(F)$ and one. Its Ramanujan sum vanishes, and
so does its induced Gauss sum. Hence every coordinate is zero.
All other orders divide $B_h$. Taking both quadratic signs may retain
extra tuples when an exponent restriction is redundant; this enlarges
the necessary set and cannot create a false exclusion.
\end{proof}

The bound $B_h$ is a bound on \emph{trace types}, not on the exponent of
an unknown category. The categorical Cauchy theorem
\cite[Theorem 3.9]{BNRW}, applied to the modular center, allows $P$ to
be the prime divisors of $|N_{K/\Q}(D)|$ (of $D$ when $K=\Q$).
Thus unbounded possible twist orders cause no omitted trace types.

\subsection{Gram columns and indicator bounds}
Fix a low matrix $L$ satisfying \eqref{eq:low}. Each remaining central
simple has a restriction vector $v=(0,v_1,\ldots,v_{r-1})$ with
nonnegative integer entries, and dimension $z_v=\sum_i v_i d_i$.
The finite set of permitted vectors contains every actual vector if we
require
\begin{equation}\label{eq:high}
 v_iv_j\le Q_{ij},\qquad Q-vv^{\mathsf T}\succeq0,\qquad
 z_v\text{ a }d\text{-number},\qquad D/z_v\text{ an algebraic integer}.
\end{equation}
Indeed, the residual Gram matrix is a sum of the rank-one matrices of
all remaining restriction vectors. The dimension conditions are the
modular-center consequences recalled in \cite[Proposition 2.1]{DP}.
Finiteness follows from $v_i\le\lfloor\sqrt{Q_{ii}}\rfloor$.
Positive semidefiniteness is checked by principal minors; for a singular
$Q$ one also imposes its rational kernel before using a nonsingular
principal block. All calculations are rational or integral.

For $j\ge3$ let
\[
 H_{i,j}=[X_i^j:1],\qquad
 \mathcal J=\{3\le j\le h:K=\Q\text{ or }F\nmid j\}.
\]
The indicator is the trace of the order-$j$ rotation on an
$H_{i,j}$-dimensional space. Its normalized relative trace to $K$ is
rational for $j\in\mathcal J$, because
$K\cap\Q(\zeta_j)=\Q$. Denote it by $\bar\nu_j(X_i)$.
If $p(j)$ is the least prime divisor of $j$, then
\begin{equation}\label{eq:indicator-bound}
 -\frac{H_{i,j}}{p(j)-1}\le \bar\nu_j(X_i)\le H_{i,j}.
\end{equation}
To see this, each normalized trace of a rotation eigenvalue is
$\mu(n)/\varphi(n)$ for some $n\mid j$. If negative, its denominator is
at least $p(j)-1$. Summing the $H_{i,j}$ terms proves the bounds.
No numerical approximation to an indicator is involved.

For every permitted vector $v$ and joint trace type $\tau$ introduce a
nonnegative real variable $x_{v,\tau}$. It relaxes the integer number of
central simples having that data. The necessary equations are
\begin{align}
 \sum_{v,\tau}x_{v,\tau}v_iv_j&=Q_{ij},\label{eq:system-gram}\\
 \ell_i+\sum_{v,\tau}x_{v,\tau}v_i z_v\tau_1&=0,
                                                       \label{eq:system-one}\\
 \ell_i+\sum_{v,\tau}x_{v,\tau}v_i z_v\tau_2&=D\epsilon_i.
                                                       \label{eq:system-two}
\end{align}
Here $\epsilon_i$ runs over every second-indicator choice allowed in
Section~\ref{sec:inputs}. For $j\in\mathcal J$, also require
\begin{equation}\label{eq:system-higher}
 \ell_i+\sum_{v,\tau}x_{v,\tau}v_i z_v\tau_j=D u_{i,j},
 \qquad u_{i,j}\in\Q,\quad
 -H_{i,j}/(p(j)-1)\le u_{i,j}\le H_{i,j}.
\end{equation}
The rationality condition is a linear equation over $\Q$: writing
$D=D_0+D_1\sqrt q$ and the left side as $w_0+w_1\sqrt q$, it is
$D_1w_0-D_0w_1=0$. The bounds can then be imposed on $w_0/D_0$,
where $D_0>0$. For $K=\Q$ the rationality equation is void.
Expanding all field coefficients, adding nonnegative slack variables,
and clearing denominators gives an integer system
\begin{equation}\label{eq:farkas}
 A x=b,\qquad x\ge0.
\end{equation}
Zero rationality rows may be retained in the rational implementation.

\begin{theorem}[Finite-indicator criterion]\label{thm:finite}
Let $R$ be a commutative fusion ring whose FP dimension field is rational
or real quadratic. Enumerate all low matrices, all second-indicator
choices and all high columns as above. If each low matrix either fails
a condition in Section~\ref{sec:inputs}, or, for every remaining indicator
choice, its system \eqref{eq:farkas} has an integer vector $y$ satisfying
\begin{equation}\label{eq:dual-ray}
 A^{\mathsf T}y\ge0,\qquad b^{\mathsf T}y<0,
\end{equation}
then $R$ has no pseudo-unitary categorification.
\end{theorem}
\begin{proof}
A hypothetical categorification supplies one of the low matrices and one
of the indicator choices. Its remaining restriction vectors satisfy
\eqref{eq:high}; Lemma~\ref{lem:trace-box} contains all their joint traces.
Equations \eqref{eq:system-gram}--\eqref{eq:system-higher} follow from
adjunction and \eqref{eq:NS}, and the indicator bounds were just proved.
Consequently it supplies a nonnegative solution $x$. But
$0\le x^{\mathsf T}A^{\mathsf T}y=b^{\mathsf T}y<0$, a contradiction.
\end{proof}

All applications use height six. Retaining higher rational-order bounds
contains the third-indicator inequality rather than introducing another
criterion. The two-moment inequalities \eqref{eq:low-moment} are kept as
cheap initial consequences. A floating-point linear program may propose
a dual ray; only the cleared integer vector and the exact inequalities
\eqref{eq:dual-ray} are accepted. Failure to find such a ray is not a
positive categorification result.

\section{Quadratic parity rigidity}\label{sec:parity}
This section removes a nonpivotal ambiguity without enumerating signed
tables. It is stated for arbitrary self-dual rank; its application to
rank five checks only tensor identities and mod-two matrices.

\begin{lemma}[Dimension-field spectral bound]\label{lem:spectral}
Let $R$ be a self-dual fusion ring, and let $\C$ be a hypothetical complex
categorification. Select the self-dual lifts in its sphericalization,
write their dimensions as $d_0=1,d_1,\ldots,d_n$, and set
\[
 K=\Q(d_1,\ldots,d_n),\qquad \delta=\sum_i d_i^2.
\]
Suppose a formal codegree $f$ of $R$ has minimal polynomial
$m_f(t)\equiv t^{\deg m_f}\pmod2$.
Then
\[
 [K:\Q]\le\operatorname{ord}_{t=0}(p_R(t)\bmod2).
\]
\end{lemma}
\begin{proof}
The lifts are self-dual: their dimensions are opposite and nonzero,
whereas spherical duality preserves dimension. Their signed quotient
has commuting real symmetric multiplication matrices $T_i$ and satisfies
\eqref{eq:signed}. It is commutative because every element of its basis
is self-dual and duality is an anti-automorphism. It is semisimple since
its real regular representation is closed under transpose.

Every formal codegree $f$ of $R$ belongs to $K$, and $\delta/f\in\OO_K$.
Indeed, inflate the corresponding character to the sphericalization.
Its formal codegree is $2f$ and the sphericalization has dimension
$2\delta$. By \cite[Theorem 2.13]{Ostrik}, a central simple has dimension
$2\delta/(2f)=\delta/f$. On forgetting, this is an integral linear
combination of the selected dimensions and their negatives. It is an
algebraic integer, and is nonzero, so $f\in K$ as well.

Put $e=[K:\Q]$. The Galois orbit of the dimension character has exactly
$e$ characters, and contributes the monic integral factor
\begin{equation}\label{eq:orbit-factor}
 Q_d(t)=\prod_{\sigma:K\hookrightarrow\CC}(t-\sigma(\delta))
\end{equation}
to the characteristic polynomial of $\sum_iT_i^2$.
Its degree is $e$, with character multiplicity, even when $\delta$
generates a proper subfield of $K$. Integrality follows from Galois
invariance and algebraic integrality of the eigenvalues; the complementary
Galois-stable factor is integral for the same reason.
In a normal overfield, reduce at any prime above two. All conjugates of
$f$ reduce to zero by the hypothesis on $m_f$. Integrality of every
conjugate of $\delta/f$ implies that all conjugates of $\delta$ do too.
Thus $Q_d(t)\equiv t^e\pmod2$.
Finally, $T_i\equiv N_i\pmod2$ gives
$t^e\mid\det(tI-\sum_iT_i^2)\bmod2=p_R(t)\bmod2$, as required.
\end{proof}

\begin{theorem}[Quadratic parity rigidity]\label{thm:parity}
Let $R$ be self-dual of rank $n+1$, with FP dimensions
\[
 s_0=1,\qquad s_i=A_i+\sqrt q\quad(1\le i\le n),
\]
where $q>1$ is squarefree, $q\equiv2$ or $3\pmod4$, and
$A_i\in\ZZ_{\ge0}$. Put $M=(N_{jj}^{\,i})_{1\le j,i\le n}$.
Suppose
\begin{equation}\label{eq:parity-hypotheses}
 1+\sum_i A_i^2=nq,\qquad
 \operatorname{ord}_{t=0}(p_R(t)\bmod2)=2,\qquad
 \ker_{\FF_2}M=\FF_2(1,\ldots,1).
\end{equation}
Assume some $A_j>0$ has irreducible fusion matrix $N_j$.
Then every complex categorification of $R$ admits a spherical structure
and has a pseudo-unitary Galois conjugate.
\end{theorem}
\begin{proof}
Use the notation of Lemma~\ref{lem:spectral}. The FP formal codegree is
\[
 D=1+\sum_i s_i^2=2nq+2\Bigl(\sum_iA_i\Bigr)\sqrt q.
\]
It generates $\Q(\sqrt q)$, and its monic minimal polynomial is
\[
 t^2-4nqt+4n^2q^2-4q\Bigl(\sum_i A_i\Bigr)^2\equiv t^2\pmod2.
\]
The lemma gives $[K:\Q]\le2$, while its proof gives $D\in K$.
Consequently $K=\Q(\sqrt q)$.

Since $\OO_K=\ZZ[\sqrt q]$, write $d_i=a_i+b_i\sqrt q$ with
$a_i,b_i\in\ZZ$. At least one $b_i$ is nonzero. Comparing irrational
coefficients in the signed square equations gives
\[
 2a_jb_j=\sum_iT_{jj}^{\,i}b_i.
\]
Let $g$ be the positive gcd of the nonzero $b_i$, and put $c_i=b_i/g$.
Divide before reducing modulo two: $Mc=0$ in $\FF_2^n$.
The integer vector $c$ is primitive, so its mod-two image is nonzero.
The kernel hypothesis forces all $c_i$ to be odd, hence all $b_i\ne0$.
This division is necessary; reducing $b$ directly could lose an even
common factor.

The characters $(1,d_i)_i$ and $(1,d_i')_i$, where prime denotes quadratic
conjugation, are distinct common eigenvectors of real symmetric matrices.
They are orthogonal, so
\begin{equation}\label{eq:orthogonality}
 1+\sum_i a_i^2=q\sum_i b_i^2.
\end{equation}
The spectral-radius inequality $|T_i|\le N_i$ bounds every conjugate
dimension by $s_i$. Therefore
\[
 |a_i|+|b_i|\sqrt q=\max(|d_i|,|d_i'|)\le A_i+\sqrt q.
\]
Since $|b_i|\ge1$, this implies $|a_i|\le A_i$. Thus
\[
 nq\le q\sum_i b_i^2=1+\sum_i a_i^2
       \le1+\sum_i A_i^2=nq.
\]
Equality holds throughout: $|b_i|=1$ and $|a_i|=A_i$.

Choose the embedding for which $|d_j|=s_j$ at an irreducible $N_j$ with
$A_j>0$, and set $v=|d|$ entrywise. From
$N_jv\ge|T_jd|=s_jv$, a positive left Perron vector forces equality.
Irreducibility then gives $v=s$, with the scalar fixed by $v_0=s_0=1$.
Put $\eta_i=\operatorname{sgn}(d_i)$, $\eta_0=1$, and
$U_{ik}^{\,\ell}=\eta_i\eta_k\eta_\ell T_{ik}^{\,\ell}$.
The vector $s$ is a positive character for both $U$ and $N$. Hence
\[
 0=\sum_\ell(N_{ik}^{\,\ell}-U_{ik}^{\,\ell})s_\ell.
\]
Each summand is nonnegative by \eqref{eq:signed}; therefore $U=N$.

Interchanging a selected lift with its opposite according to $\eta_i$
realizes this integer sign change. The selected and opposite
multiplicities in a product are $(N+U)/2$ and $(N-U)/2$. The latter
vanish, so the selected lifts form a fusion subcategory of the
sphericalization, equivalent to $\C$ under the forgetful functor.
The spherical structure restricts to it; compare
\cite[Section 5.1]{Ostrik}. Although one embedding was used in the
argument, $U=N$ is an equality of integer tensors and constructs the
section for the original category. Applying the chosen embedding to
its structure data gives positive dimensions $s_i$, hence the stated
pseudo-unitary Galois conjugate.
\end{proof}

\begin{remark}[Recognition from the tensor]
In rank five the vector $(A_i)$ need not be guessed. The irrational
coefficient of $s_j^2=\sum_iN_{jj}^i s_i$ gives
$2A_j=\sum_{i=1}^4N_{jj}^i$. Recover these four integers from the square
rows, set $q=(1+\sum_iA_i^2)/4$, and verify every dimension-product
identity. Squarefreeness, the mod-two characteristic-polynomial
multiplicity, the kernel and irreducibility are then exact finite tests.
The supplement implements this recognition for every census tensor;
no radicand or list of exceptional indices is supplied to the criterion.
\end{remark}

\section{Cyclic order-four central support}\label{sec:cyclic}
Write $K(C_4,4k)$ for the near-group ring
\[
 g^4=1,\qquad gX=Xg=X,\qquad
 X^2=1+g+g^2+g^3+4kX,
\]
where $k\ge1$ is an integer. The following reduction is uniform in $k$.
Thornton's theorem \cite[Theorem 1.3]{Thornton} permits a pseudo-unitary
Galois conjugate of any complex categorification of a near-group ring;
no braidedness or unitarity hypothesis is imposed on that reduction.
Set
\[
 d=2k+2\sqrt{k^2+1},\quad d^2=4+4kd,\quad
 D=4+d^2=8+4kd,\quad t=D/D'=1+kd=d^2/4.
\]
The dimension $d$ is irrational because $k^2<k^2+1<(k+1)^2$.

\subsection{A finite central-support lemma}
For a central simple $Z$ put
\[
 F(Z)=\sum_{r=0}^3m_{rZ}g^r+b_ZX,\qquad a_Z=\sum_r m_{rZ}.
\]
Call $Z$ \emph{mixed} if $a_Z>0$, and \emph{pure} otherwise.
The induction formula and adjunction give
\begin{equation}\label{eq:ng-gram}
 \sum_Zm_{rZ}m_{sZ}=4\delta_{rs}+1,\qquad
 \sum_Zm_{rZ}b_Z=4k,\qquad
 \sum_Zb_Z^2=8+16k^2.
\end{equation}

\begin{lemma}\label{lem:ng-support}
The mixed simples consist of six objects $B_{rs}$, $0\le r<s\le3$,
and two objects $E_{r,1},E_{r,2}$ for every $r$. Their restrictions are
\[
 F(B_{rs})=g^r+g^s+kX,\qquad
 F(E_{r,1})=g^r+u_rX,\quad F(E_{r,2})=g^r+(k-u_r)X,
\]
where $u_0=0$, $0\le u_r\le\lfloor k/2\rfloor$, and $u_1=u_3$.
For some $\eta\in\{i,-i\}$ their twists are
\begin{equation}\label{eq:ng-twists}
\begin{array}{c|c}
 E_{0,*},E_{2,*},B_{01},B_{02},B_{03}&1\\
 B_{12},B_{23}&-1\\
 B_{13}&\eta\\
 E_{1,*},E_{3,*}&-\eta.
\end{array}
\end{equation}
The set of $r$ for which $u_r=0$ is a subgroup of $C_4$. Its size
$h$ belongs to $\{1,2,4\}$ and equals the number of central invertibles.
Writing $L_r=u_r^2+(k-u_r)^2$, the total pure mass is
\begin{equation}\label{eq:ng-pure}
 \sum_{Z\,\mathrm{pure}}b_Z^2=8+10k^2-\sum_rL_r.
\end{equation}
If $h\ge2$, there is a central invertible $z$ with $F(z)=g^2$, $z^2=1$
and $\theta_z=1$. Splitting pure simples by their monodromy sign with
$z$ gives
\begin{align}
 \sum_{Z\in\mathcal P_+}b_Z^2&=4+6k^2-\sum_rL_r,
 &\sum_{Z\in\mathcal P_-}b_Z^2&=4+4k^2,\label{eq:ng-split}\\
 \sum_{Z\in\mathcal P_-}b_Z^2\theta_Z^j&=0
 &&\text{for every odd }j.\label{eq:ng-odd}
\end{align}
\end{lemma}
\begin{proof}
The formal codegrees are $D,D',4,4,4$. Thus
$I(1)=1\oplus A\oplus B_{01}\oplus B_{02}\oplus B_{03}$, with
$F(A)=1+kX$ and $F(B_{0r})=1+g^r+kX$; this follows from their dimensions
$1,t,D/4,D/4,D/4$ and \eqref{eq:ng-gram}. All have twist one.

The pointed associator class is trivial. Indeed, $\operatorname{add}(X)$
is a rank-one module category over $\C_{\mathrm{pt}}$; choosing the
isomorphisms $g^rX\simeq X$ gives a two-cochain whose coboundary is that
associator. Thus the pointed indicators equal those of $\Vec_{C_4}$:
$\nu_j(g^r)=1$ if $g^{rj}=1$ and zero otherwise. Positivity in
\eqref{eq:NS} at $j=\operatorname{ord}(g^r)$ forces
$m_{rZ}>0\Rightarrow\theta_Z^{\operatorname{ord}(g^r)}=1$.
All mixed twists therefore belong to $\{1,-1,i,-i\}$.

Apply the four-point Fourier transform to the traces for $I(g^r)$ at
$j=0,1,2,3$, with $j=0$ denoting dimension. Comparing coefficients of
$1,d$ is valid over $\Q(i)$ because $d\notin\Q(i)$. After summing over
$r$ the budgets are
\begin{equation}\label{eq:ng-budgets}
\begin{array}{c|rrrr}
 \theta&1&-1&i&-i\\\hline
 \sum a_Z^2&16&8&4&4\\
 \sum a_Zb_Z&8k&4k&2k&2k.
\end{array}
\end{equation}
Every $m_{rZ}$ is zero or one. For $r=0$ this follows from $I(1)$.
For $r\ne0$, $B_{0r}$ already contributes one to the diagonal Gram
entry five. An additional multiplicity two would consume the remaining
four; its overlap with either of the other two nonzero rows would then
be even, contrary to the off-diagonal entry one in \eqref{eq:ng-gram}.

A support of size three must avoid the unit row and therefore consist
of $g,g^2,g^3$. Its twist is $\pm1$ and it contributes nine to its
$a^2$-budget. But the negative budget is eight, while $I(1)$ already
contributes fourteen to the positive budget sixteen. A support of size
four is already impossible in the unit row. Each pair of rows
consequently overlaps in exactly one pair-support object, and each row
has two remaining single-support objects.

Neither $B_{12}$ nor $B_{23}$ fits in the remaining positive budget two,
so both have twist $-1$ and exhaust its budget. The two $g^2$-supported
objects must then have twist one. The object $B_{13}$ exhausts one
imaginary budget, leaving all four $g$- and $g^3$-supported singles with
the opposite imaginary twist. This proves \eqref{eq:ng-twists}.
The individual Fourier equations for $I(g)$ and $I(g^3)$ give mass
$D/4=2+kd$ in each fourth-root sector; those for $I(g^2)$ give $D/2$
in each sign sector. Comparing dimensions forces every pair coefficient
to be $k$ and every single pair to have sum $k$. Duality identifies the
unordered coefficient pairs over $g$ and $g^3$, so $u_1=u_3$.

Only a single-support object with coefficient zero has dimension one.
Forgetting embeds the group of central invertibles into $C_4$, because
the only invertible summand of $I(1)$ is the unit. This proves the
subgroup assertion and \eqref{eq:ng-pure}. If $h\ge2$, its order-two
object lies over $g^2$ and has twist one by \eqref{eq:ng-twists}.
Tensoring with it translates supports by two. All eight singles and
exactly $B_{02},B_{13}$ centralize it; the other four pairs have monodromy
$-1$. Their total mixed squared dimension is
\[
 16+16kd+(\sum_rL_r+2k^2)d^2.
\]
The centralizer has dimension $D^2/2$ by
\cite[Theorem 3.14]{DGNO}, and
\[
 D^2/2=16+16kd+(4+8k^2)d^2.
\]
Subtracting proves \eqref{eq:ng-split}. On the negative sector tensoring
with $z$ preserves $b_Z$ and changes the twist to $-\theta_Z$.
There are no fixed points, and pairing the terms proves
\eqref{eq:ng-odd}.
\end{proof}

\begin{lemma}[Balanced Galois obstruction]\label{lem:ng-galois}
A branch with $h=1$ and $u_1=u_2=u_3=k/2$ requires $k^2/8\in\ZZ$.
\end{lemma}
\begin{proof}
Here $A$ is the unique central simple of dimension $t$.
The signed Galois symmetry of the normalized modular matrix $S/D$
\cite[(2.13)]{BNRW}, at the embedding sending $d$ to $d'$, sends the
unit row to the $A$ row with positive sign. Hence
$S_{AZ}=t\,\sigma(\dim Z)$. Its eigenvalue ratios
$S_{AZ}/\dim Z$ are $t,t^{-1}$ at $1,A$, one at every other mixed
simple, and minus one at every pure simple; use $td'=-d$.
Since $A$ is self-dual, subtracting the unit fusion coefficient from
Verlinde's formula leaves only the two exceptional columns:
\[
 N_{AA}^Z-\delta_{Z,1}
 =\frac{(t^2-1)(\dim Z-\sigma(\dim Z))}{D^2}
 =\frac{k}{4}b_Z.
\]
The last equality uses $(t^2-1)(d-d')/D^2=k/4$.
Every single over a nonunit group element has $b_Z=k/2$.
Its fusion coefficient must therefore equal the integer $k^2/8$.
\end{proof}

\subsection{The parameterized indicator system}
Fix a moment height $H\ge3$. For each $u_1,u_2,u_3$ in Lemma~\ref{lem:ng-support} and
$\epsilon=\nu_2(X)\in\{-1,1\}$, take the known mixed data from
\eqref{eq:ng-twists}. Put
\[
 W_j=\sum_{Z\,\mathrm{mixed}} b_Z(a_Z+b_Zd)\theta_Z^j.
\]
Apply normalized relative trace to $K=\Q(\sqrt{k^2+1})$, interpreted
using its squarefree radicand. We may fix $\eta=i$, since the resulting
traces are unchanged by $i\mapsto-i$. Assign a nonnegative weight
$x_\tau$ to each pure joint trace type of Lemma~\ref{lem:trace-box}; it
represents the sum of $b_Z^2$ for that type. Its total is
\eqref{eq:ng-pure}. For $h\ge2$, instead use two pools with the masses
\eqref{eq:ng-split}. In the negative pool replace a joint tuple by
\[
 (\tau_j)_{j=1}^H\longmapsto
 \bigl(0\text{ for odd }j,\ \tau_j\text{ for even }j\bigr)_{j=1}^H,
\]
which is the average of the paired twists $\theta,-\theta$.
The first two equations are
\[
 \Tr^{\mathrm{norm}}_{L/K}(W_1)+d\sum x_\tau\tau_1=0,\qquad
 \Tr^{\mathrm{norm}}_{L/K}(W_2)+d\sum x_\tau\tau_2=\epsilon D.
\]
For orders $3\le j\le H$ with $F\nmid j$, impose the rationality and
bounds \eqref{eq:indicator-bound} on
$(\Tr^{\mathrm{norm}}_{L/K}(W_j)+d\sum x_\tau\tau_j)/D$.
The needed rotation dimensions are obtained without tensor powers:
if $X^j=\alpha_j\sum_rg^r+\beta_jX$, then
\[
 (\alpha_1,\beta_1)=(0,1),\qquad
 (\alpha_{j+1},\beta_{j+1})=(\beta_j,4\alpha_j+4k\beta_j),\qquad
 [X^j:1]=\alpha_j.
\]
The prime support is that of $2(k^2+1)$, since
$N_{K/\Q}(D)=64(k^2+1)$. All these systems are rational and become
integer systems after row scaling.

\begin{theorem}[Cyclic-four central-support criterion]\label{thm:cyclic}
Fix $k\ge1$ and a finite moment height $H\ge3$. Enumerate all the
coefficient and indicator branches above. If each branch is rejected by
Lemma~\ref{lem:ng-galois} or by an integer separating vector
\eqref{eq:dual-ray} for its central-support system, then $K(C_4,4k)$
admits no fusion categorification over $\CC$.
\end{theorem}
\begin{proof}
Any categorification has a pseudo-unitary Galois conjugate. Its central
support satisfies Lemma~\ref{lem:ng-support}, so belongs to an enumerated
branch. Lemma~\ref{lem:ng-galois} rules out the indicated balanced
branches. Every other branch supplies nonnegative weights in its
indicator system by \eqref{eq:NS}, the paired negative sector and
Lemma~\ref{lem:trace-box}. An integer separator contradicts their
existence exactly as in Theorem~\ref{thm:finite}.
\end{proof}

The software recognizes the entire cyclic-four tensor before applying
this criterion; recognizing only the number of invertible objects is
not sufficient. At height six it rejects the parameters $k=2,3,4,5$
encountered in the census beyond the positive case $k=1$. Their allocations are
reported together in Section~\ref{sec:automatic}; no separate numerical
root-of-unity argument is used for one selected ring.

\section{Unitary realizations}\label{sec:models}
\begin{table}[ht]\centering\small
\begin{tabular}{rrl}\toprule
$(\mu,j)$&Local index&Unitary realization\\\midrule
$(1,1)$&7&Tambara--Yamagami for $C_2^2$\\
$(1,3)$&9&$SU(2)_4$\\
$(1,4)$&14&$\Rep(D_{14})$\\
$(1,6)$&12&$\Rep(S_4)$\\
$(1,7)$&11&$PSU(2)_8$\\
$(1,10)$&15&$PSU(2)_9$\\
$(1,11)$&2&Tambara--Yamagami for $C_4$\\
$(1,13)$&4&The dimension-twelve model of \cite{LPR}\\
$(1,14)$&5&The dimension-twenty-four model of \cite{LPR}\\
$(1,16)$&1&$\Vec(C_5)$\\
$(2,27)$&33&$\Rep(A_5)$\\
$(2,36)$&5&$PSU(3)_4$\\
$(2,37)$&1&$\Rep(C_7\rtimes C_3)$\\
$(3,40)$&28&$C_2$-equivariantization of an Izumi $C_3+6$ category\\
$(3,75)$&8&$\Rep(\operatorname{AGL}(1,5))$\\
$(4,55)$&20&Izumi near-group $C_2^2+4$\\
$(4,123)$&6&Izumi near-group $C_4+4$\\
\bottomrule\end{tabular}
\caption{The seventeen positive rings. $D_{14}$ has order fourteen. Each tensor is constructed and matched by a unit-fixing permutation; the supplied permutation is checked.}\label{tab:models}\end{table}

The pointed, Tambara--Yamagami, finite-group and quantum-group entries
use their standard unitary realizations. The two nonstandard
multiplicity-one entries are the tensors of \cite[Section 3.5]{LPR},
with the unitary constructions in its Sections 4.2 and 4.4. The
$PSU(3)_4$ tensor agrees with the modular data in \cite{HR}.
The near-group realizations for $C_4+4$ and $C_2^2+4$ are
\cite[Examples 9.5--9.6]{Izumi}. These references establish existence,
not negative classifications used as an oracle.

The software constructs the $SU(2)$ tensors by the truncated
Clebsch--Gordan rule, the finite-group tensors by exact character
products, and the remaining displayed tensors in the cited conventions.
It checks their ring axioms, distinctness and a full tensor isomorphism
to the pinned census tensors. The unitary existence is supplied by the
constructions, not inferred just from passing the ring axioms.

For completeness, we explain the equivariantization entry, since an
order-two outer symmetry by itself would not supply the required action.
In \cite[Example 13.10]{Izumi}, the $C_3+6$ Cuntz realization has actual
symmetries generated by the pairs $(-1,R_{\pi/4})$ and
$(\id,L)$, where $L=\operatorname{diag}(1,-1)$ and $R_{\pi/4}$ is planar
rotation. Their product is $(-1,U)$ with $U=R_{\pi/4}L$ and $U^2=I$.
Thus it gives an actual order-two automorphism of the realization,
fixing the noninvertible endomorphism and the unit and interchanging
the two nontrivial invertibles. It fixes the unit inclusion in the square
of that endomorphism. On its six-dimensional self-fusion multiplicity
space $\CC[C_3]\otimes\CC^2$ it is inversion tensored with the
reflection $U$, and has trace zero. The two eigenspaces have dimension
three.

The $C_2$-equivariantization has simples $1,g,V,X,Y$:
two equivariant structures over the unit and the noninvertible object,
and one object over the orbit of the two other invertibles. The fixed
unit inclusion makes $X,Y$ self-dual. Tensoring equivariant structures
and using the two three-dimensional eigenspaces gives
\begin{align*}
 g^2&=1,&gV&=V,&gX&=Y,\\
 V^2&=1+g+V,&VX&=VY=X+Y,\\
 X^2=Y^2&=1+V+3X+3Y,&XY&=g+V+3X+3Y.
\end{align*}
The action is unitary, so its equivariantization is unitary. Its tensor
is generated from these products and matched exactly, just like the
other sixteen models.

\section{Automatic application and the next frontier}\label{sec:automatic}
\subsection{A label-free decision procedure}
Every exclusion receives its input as the complete integer multiplication
tensor. The application order is arithmetic screening, the cyclic-four
central-support criterion when recognized, quadratic parity when its
hypotheses hold, and otherwise the positive-indicator sphericalization
framework. Pseudo-unitary branches use the low-induction bounds followed
by the same finite-indicator systems. Spherical branches exhaust the
dimension-character components. Signed branches exhaust the bounded
associative tables in the normalization \eqref{eq:positive-lifts} and
every eligible character component. Successful branches are rechecked
before a certificate is accepted.

None of these functions receives a multiplicity index, historical label,
previous status or positive-list complement as a hypothesis. The
seventeen model matches are a separate positive channel. An arithmetic
failure, field obstruction, missing complete low matrix, exact moment
violation or integer separating vector is a mathematical witness. A
search timeout, unsupported field or unsuccessful linear program is not.
This policy is enforced by the decision ledger.

The top-level allocations in Table~\ref{tab:allocation} are disjoint and
follow the actual decision order. Formal-codegree arithmetic comes first
because it settles almost all inputs without a central search. The
cheap cyclotomic and fusion-subring checks follow. Exact recognition of
cyclic-four and quadratic-parity tensors is cheap; their specialized
systems are built only on those families. Positive reduction and central
induction precede the full signed sphericalization search. These choices
balance rejection frequency and cost, without claiming a formal runtime
optimality theorem. The low-moment bounds are retained for their cost
advantage, despite being consequences of stronger moment conditions.
\begin{table}[ht]\centering\small
\begin{tabular}{lrrrr}\toprule
First accepted top-level certificate&First $\mu$&$\mu\le17$&$\mu\le18$&$\mu\le20$\\\midrule
Formal-codegree $d$-number test&1&23075&28124&40558\\
Cyclotomicity test&5&5&6&8\\
Fusion-subring index test&3&12&12&14\\
Cyclic-four central support&8&3&3&4\\
Quadratic parity and positive induction&3&12&12&20\\
Quadratic parity and spherical characters&12&6&6&6\\
Positive reduction and central induction&2&27&29&30\\
Positive-indicator sphericalization&2&256&282&340\\
\midrule Total general exclusions&&23396&28474&40980\\
\bottomrule\end{tabular}
\caption{Disjoint first-accepted certificates in executable order. First use at a multiplicity does not date the underlying theorem. The order favors cheap general tests, then structural recognition before restricted and full central-induction searches.}\label{tab:allocation}\end{table}

\begin{table}[ht]\centering\small
\setlength{\tabcolsep}{3pt}
\begin{tabular}{r|r|rrrrrrrr|r}\toprule
$\mu$&To decide&D&Cy&Sub&C4&QP&QS&Gate&Sign&Open\\\midrule
1&6&6&0&0&0&0&0&0&0&0\\
2&34&26&0&0&0&0&0&4&4&0\\
3&80&62&0&2&0&2&0&6&8&0\\
4&132&119&0&0&0&0&0&2&11&0\\
5&209&197&2&1&0&0&0&0&9&0\\
6&336&312&0&1&0&0&0&4&19&0\\
7&477&461&0&1&0&0&0&0&15&0\\
8&733&708&0&0&1&0&0&1&23&0\\
9&863&848&1&0&0&0&0&1&13&0\\
10&1082&1064&0&0&0&1&0&2&15&0\\
11&1383&1365&0&1&0&0&0&2&15&0\\
12&1948&1920&0&2&1&1&1&1&22&0\\
13&2211&2188&0&0&0&2&3&0&18&0\\
14&2554&2531&1&1&0&0&1&3&17&0\\
15&2936&2909&1&2&0&1&1&0&22&0\\
16&4157&4125&0&1&1&0&0&1&29&0\\
17&4255&4234&0&0&0&5&0&0&16&0\\
18&5080&5049&1&0&0&0&0&2&26&2\\
19&5640&5609&1&2&0&3&0&0&25&0\\
20&6866&6825&1&0&1&5&0&1&33&0\\
\bottomrule\end{tabular}
\caption{Exact-multiplicity exclusion sequence. To decide omits independently constructed unitary models. Stage symbols follow the row order in Table~\ref{tab:allocation}; each ring appears once, at its first accepted certificate. Open concerns general categorification by these criteria.}\label{tab:stages}\end{table}

\begin{table}[ht]\centering\small
\begin{tabular}{rrrrr}\toprule
$\mu$&Total&Unitary models&General exclusions&General residuals\\\midrule
1&16&10&6&0\\
2&37&3&34&0\\
3&82&2&80&0\\
4&134&2&132&0\\
5&209&0&209&0\\
6&336&0&336&0\\
7&477&0&477&0\\
8&733&0&733&0\\
9&863&0&863&0\\
10&1082&0&1082&0\\
11&1383&0&1383&0\\
12&1948&0&1948&0\\
13&2211&0&2211&0\\
14&2554&0&2554&0\\
15&2936&0&2936&0\\
16&4157&0&4157&0\\
17&4255&0&4255&0\\
18&5080&0&5080&0\\
19&5640&0&5640&0\\
20&6866&0&6866&0\\
\bottomrule\end{tabular}
\caption{Exact automatic coverage of all pinned rank-five tensors through multiplicity twenty. All general, pivotal, spherical, pseudo-unitary and unitary existence questions in this range are decided.}\label{tab:coverage}\end{table}


The last new \emph{criterion family} first used in the checked range is
cyclic-four central support at multiplicity eight: from multiplicity nine
through twenty, every exclusion uses the same repertoire. The
parity-plus-spherical-character \emph{certificate type}, which combines
already available conditions, first occurs at multiplicity twelve.
Thus every top-level certificate type has already appeared by twelve,
and from thirteen through twenty the executable proof allocation
introduces no new type. This identifies the earliest multiplicity in
the verified ledger after which no new criterion family, respectively
certificate type, is needed. It makes no claim of a uniform unbounded
classification: two general residuals occur at eighteen.

For the cyclic-four tensors, automatic recognition returns $k=2,3,4,5$.
The height-six criterion checks respectively $6,6,14,14$ support/indicator
branches. At $k=2$ the balanced Galois test removes two and exact
separators remove four; at $k=3,4,5$ all branches have exact separators.
The same routine leaves the positive $k=1$ model unexcluded.
These are recorded applications of Theorem~\ref{thm:cyclic}, not three
separate nonexistence arguments. Similarly, the twenty-six top-level
quadratic-parity applications are detected from the square rows and
verified modulo two; no fixed list of radicands enters the code.

Integral inputs use the rational specialization of
Theorem~\ref{thm:finite}. No published classification of integral rings
is invoked to exclude them. No one-ring two-moment calculation, special
root-length test or isolated signed-case file is needed. All signed
certificates use the positive-indicator normalization, rather than
retaining separate ordinary and adjoint-parity enumeration modes in the
proof allocation.

\subsection{Proof of Theorem~\ref{thm:main}}
The pinned repository census \cite{Census} gives the counts in
Table~\ref{tab:census}. Section~\ref{sec:models} supplies the
seventeen unitary realizations and their exact tensor matches.
For every other ring with $\mu\le17$, the exhaustive ledger supplies a
general exclusion certificate covered by the theorems and cited inputs
above. There are $23396$ such certificates. At multiplicity eighteen,
$5078$ of the $5080$ rings have general exclusion certificates. At exact
multiplicities nineteen and twenty, every one of the $5640+6866$
repository tensors has a general exclusion certificate. These later
certificates use the same criterion order and exact checker, but are
addressed directly against the pinned repository tensor order, not the
older local basis labels. The two remaining multiplicity-eighteen
rings have complete spherical exclusions. Their bases are
self-dual, so every pivotal structure would be spherical: the left
pivotal dimension of a simple equals the right dimension of its dual.
This excludes pivotal existence as well. The inclusions among the five
existence notions and the seventeen unitary realizations give all the
asserted equivalences and ranges.\qed

The exhaustive step here is a finite, exact certificate verification;
it is not a proof-assistant formalization of the categorical theorems.
The release includes every tensor, branch and integer witness, so its
finite assertions can be reproduced from the pinned census rather than
trusted as a numerical census label. The two proof-ledger verification
commands recheck the entire $28493+12506=40999$-ring source and each
accepted certificate.

\subsection{The two general residuals}
In the twenty self-dual coordinates of Section~\ref{sec:data}, the two
remaining tensors are
\begin{equation}\label{eq:frontier-tensors}\begin{aligned}
v^{(5073)}&=(16,16,17,17,17,17,17,15,17,17,\\ &\hspace{15mm}17,17,17,16,17,18,18,17,17,18),\\
v^{(5076)}&=(17,16,16,17,16,17,17,17,17,17,\\ &\hspace{15mm}17,16,17,17,17,17,17,17,18,18).
\end{aligned}\end{equation}

They have common sorted FP type
\[
 (1,a,a,a+1,a+2),\qquad a=33+2\sqrt{285},
 \qquad D=9120+540\sqrt{285}.
\]
Their formal-codegree polynomials are respectively
\begin{align*}
 (t-3)^2(t-15)(t^2-18240t+68400),\\
 (t-3)(t-5)^2(t^2-18240t+68400).
\end{align*}
Both reduce to $t^2(t+1)^3$ modulo two. Therefore
Lemma~\ref{lem:spectral} restricts a hypothetical sphericalization's
entire dimension field to $\Q(\sqrt{285})$, not merely the field of its
global dimension. However, the displayed FP dimensions do not satisfy
Theorem~\ref{thm:parity}: $285\equiv1\pmod4$ and the irrational
coefficient is two. The applicable ring of integers is
$\ZZ[(1+\sqrt{285})/2]$. Extending the parity rigidity argument to this
situation is a specific possible continuation, not a theorem used here.

At both tensors, a complete spherical certificate was obtained; the
general positive-indicator signed enumeration reached its sixty-second
subprocess limit without completion. The discovery records preserve this
distinction. Thus multiplicity eighteen is the first unresolved
\emph{general} stage of the present pipeline after complete treatment of
sixteen and seventeen. Neither tensor is certified feasible, and neither
is asserted to be unknown throughout the literature.

\subsection{Numbering and reproducibility}
The release's notation $\F{m}{j}$ preserves the first index and basis
from the historical ancillary file of \cite{VS} when one exists.
Classes absent from that file have appended archival labels in local
tensor order. Separately, the field \texttt{local\_index} gives the
one-based ordinal within the archived tensor records at fixed $m$.
At multiplicity eighteen the two indices coincide. Fingerprints
of the full tensors, rather than these labels, key every certificate.

The command \texttt{python verify.py} checks the first $28493$ tensors,
canonical uniqueness, the full pinned-repository isomorphism crosswalk,
and all $28474$ general exclusion
certificates, the seventeen constructed model matches and both spherical
frontier certificates. The companion command
\texttt{python software/verify\_extension.py} checks exact source order,
ring axioms, uniqueness and all $12506$ later general exclusion
certificates. Both regenerate each finite matrix and signed branch.
Python/native Gram comparisons and deliberate certificate mutations are
additional controls. The command \texttt{python -O verify.py} performs
the same proof checks: none is implemented using a removable Python
\texttt{assert}. The accepted separating vectors are tested with integers;
a numerical optimization solver is used only in discovery.

The verification guide gives exact commands, input conventions, checked
integer ranges, dependencies and the failure policy. The bundled pinned
repository tensors are the census input. The discovery command
accepts either an external tensor or a canonical index and
applies the same criteria at the chosen moment height and search budget.
The archive also contains complete decision and provenance ledgers,
input hashes and provenance notes.

\section*{AI use disclosure}
GPT-6 Astra assisted with mathematical exploration, code development, and manuscript preparation.
\begin{thebibliography}{99}
\raggedright
\small
\bibitem{ABDP} M. A. Alekseyev, W. Bruns, J. Dong and S. Palcoux,
\emph{Classifying integral Grothendieck rings up to rank 5 and beyond},
accepted for publication in \emph{Letters in Mathematical Physics};
\href{https://arxiv.org/abs/2507.07023}{arXiv:2507.07023}.

\bibitem{BNRW} P. Bruillard, S.-H. Ng, E. C. Rowell and Z. Wang,
\emph{Rank-finiteness for modular categories}, J. Amer. Math. Soc.
\textbf{29} (2016), 857--881;
\href{https://arxiv.org/abs/1310.7050v3}{arXiv:1310.7050v3}.
\bibitem{DP} \emph{Rank-four Grothendieck rings through multiplicity 128},
manuscript and exact-certificate supplement (2026).
The numbered results cited here refer to the finite-classification manuscript.
\bibitem{Census} \emph{Fusion-ring census},
\href{https://github.com/sebastienpalcoux/fusion-ring-census/tree/b86eb6dc6eeeafe5fc5f6906d7b074e0322e9726}{repository at commit \texttt{b86eb6dc6eeeafe5fc5f6906d7b074e0322e9726}},
\texttt{results/rank5/counts.csv} and
\texttt{results/rank5/tables.txt.gz} (2026).
\bibitem{HR} S.-M. Hong and E. C. Rowell,
\emph{On the classification of the Grothendieck rings of non-self-dual
modular categories}, J. Algebra \textbf{324} (2010), 1000--1015;
\href{https://arxiv.org/abs/0907.1051}{arXiv:0907.1051}.
\bibitem{Izumi} M. Izumi,
\emph{A Cuntz algebra approach to the classification of near-group categories},
in \emph{Proceedings of the 2014 Maui and 2015 Qinhuangdao Conferences
in Honour of Vaughan F. R. Jones' 60th Birthday},
Proc. Centre Math. Appl. Austral. Nat. Univ. \textbf{46} (2017), 222--343;
\href{https://arxiv.org/abs/1512.04288v2}{arXiv:1512.04288v2}.
\bibitem{NS} S.-H. Ng and P. Schauenburg,
\emph{Frobenius--Schur indicators and exponents of spherical categories},
Adv. Math. \textbf{211} (2007), 34--71;
\href{https://arxiv.org/abs/math/0601012}{arXiv:math/0601012}.
\bibitem{LPR} Z. Liu, S. Palcoux and Y. Ren,
\emph{Classification of Grothendieck rings of complex fusion categories
of multiplicity one up to rank six}, Lett. Math. Phys. \textbf{112}
(2022), article 54;
\href{https://arxiv.org/abs/2010.10264v5}{arXiv:2010.10264v5}.
\bibitem{VS} G. Vercleyen and J. K. Slingerland,
\emph{On low rank fusion rings}, J. Math. Phys. \textbf{64} (2023),
article 091703;
\href{https://arxiv.org/abs/2205.15637v4}{arXiv:2205.15637v4}.
Input: the ancillary file \texttt{FusionRingMultiplicationTables}.

\bibitem{DGNO} V. Drinfeld, S. Gelaki, D. Nikshych and V. Ostrik,
\emph{On braided fusion categories I}, Selecta Math. (N.S.)
\textbf{16} (2010), 1--119;
\href{https://arxiv.org/abs/0906.0620v3}{arXiv:0906.0620v3}.
\bibitem{Ostrik} V. Ostrik,
\emph{Pivotal fusion categories of rank 3}, Mosc. Math. J.
\textbf{15} (2015), no.~2, 373--396;
\href{https://arxiv.org/abs/1309.4822v1}{arXiv:1309.4822v1}.
\bibitem{Thornton} J. Thornton, \emph{On braided near-group categories},
\href{https://arxiv.org/abs/1102.4640v1}{arXiv:1102.4640v1} (2011).
\end{thebibliography}

\end{document}
